% =============================================================================
\section{Tarefa 4 --- Orientação com o giroscópio}
\label{sec:gyro}
% =============================================================================

A orientação obtém-se integrando a velocidade angular corrigida. Para evitar
a singularidade dos ângulos de Euler (a \ang{90} de arfagem dois eixos
alinham-se, como na singularidade do pulso do robô no Laboratório~2), a
orientação é guardada num quaternião unitário $q$ e integrada
por~\cite{titterton2004}
\begin{equation}
\dot q = \tfrac12\, q \otimes (0;\ \vect\omega), \qquad
q_{k+1} = q_k \otimes \exp\!\bigl(\tfrac12\,\vect\omega_k\,\Delta t\bigr),
\label{eq:quat}
\end{equation}
com normalização a cada passo. Os ângulos só são calculados no fim, para
apresentação. A orientação inicial vem do inclinómetro, com o sensor parado.

O método responde depressa e não é afetado pela aceleração do sensor, mas o
erro cresce com o tempo: um desvio residual $\delta b_g$ acumula
$\delta b_g\,t$, e o ruído soma-se como um passeio aleatório.

\todo{resultados no ensaio D (erro ao longo do tempo) e no E (erro ao
regressar à pose inicial; deriva em \si{\degree\per\minute}), com e sem a
calibração da Tarefa~2.}
